\chapter{Implementation and Results}
\label{ch:impl}

\section{Software Architecture}
\label{sec:arch}

The system is implemented in Python without a robotics middleware layer.
Avoiding ROS was a deliberate choice: the pipeline is a batch data-processing
system rather than a distributed real-time control system, and the node graph,
message serialisation and build tooling that middleware provides would add
operational complexity without addressing any requirement of the offline
setting. Communication with the manipulator uses \texttt{ur-rtde} directly.

The implementation is organised into the layers listed in
Table~\ref{tab:modules}. The separation follows the pipeline stages of
Fig.~\ref{fig:pipeline}, with the coordinate-frame contract---every spatial
quantity carries an explicit frame identifier---enforced at module boundaries so
that frame errors surface as type errors rather than as silent numerical
mistakes.

\begin{table}[htbp]
\caption{Principal Modules of the Implementation}
\label{tab:modules}
\centering
\small
\begin{tabular}{p{32mm}p{45mm}p{45mm}}
\toprule
Module & Responsibility & Principal dependencies \\
\midrule
\texttt{capture} & Device drivers, hardware sync, host-clock stamping &
  Azure Kinect SDK, \texttt{pyrealsense2} \\
\texttt{calibration} & ChArUco detection, intrinsics, extrinsics, robot
  registration & OpenCV \\
\texttt{perception} & Landmark detection, depth lifting, confidence-weighted
  fusion & MediaPipe, OpenCV \\
\texttt{representation} & Object pose, object-relative transform, spline
  interpolation & NumPy, SciPy \\
\texttt{retargeting} & Cost functional assembly and solution & Pink, Pinocchio \\
\texttt{replay} & Trajectory execution, proprioception logging, screening &
  \texttt{ur-rtde} \\
\texttt{export} & LeRobot/HDF5 serialisation & h5py, LeRobot \\
\texttt{interface} & Recording control and inspection & FastAPI, Streamlit \\
\bottomrule
\end{tabular}
\end{table}

The Azure Kinect depth engine requires a software rendering context inside the
containerised deployment; this is obtained by forcing software rasterisation
through the appropriate driver environment variables. The system is released
under the MIT licence.

\section{Experimental Setup}
\label{sec:setup}

\subsection{Hardware configuration}

Experiments used a UR3 manipulator with a parallel-jaw gripper, two Azure
Kinect DK cameras in master/subordinate hardware synchronisation and one Intel
RealSense D435i, all observing a tabletop workspace anchored by a ChArUco
board. Camera placement was near-orthogonal to preserve view overlap
(Section~\ref{sec:hardware}). Recording ran at \ph{XX}~Hz.

\subsection{Task suite}
\label{sec:tasks}

Evaluation used a deliberately constrained task class, chosen so that the
quasi-static and rigid-object assumptions of Section~\ref{sec:assumptions} hold
and so that success is unambiguously measurable:

\begin{enumerate}[leftmargin=2em]
  \item \emph{Pick-and-place.} Grasp a rigid object and place it at a marked
        target location.
  \item \emph{Pick-and-place with displacement.} As above, but with the object
        initialised at a pose not seen during demonstration, to test the
        object-relative representation.
  \item \emph{Pouring.} Grasp a container and tip it over a target, testing
        orientation fidelity rather than position alone.
  \item \ph{[Additional task, if included.]}
\end{enumerate}

For each task, \ph{N} demonstrations were recorded by \ph{one} demonstrator.

\subsection{Baselines and ablations}

The proposed formulation \eqref{eq:cost} is compared against:

\begin{enumerate}[leftmargin=2em]
  \item \textbf{B1 --- Sequential.} Task-space low-pass filtering, per-frame
        inverse kinematics, joint-space low-pass filtering. This is the
        pattern used by offline pipelines such as \cite{okami} and the
        principal comparison for H1.
  \item \textbf{B2 --- Joint, no confidence.} Equation~\eqref{eq:cost} with
        $\omega_t \equiv 1$. Isolates the contribution of confidence weighting
        for H2.
  \item \textbf{B3 --- Joint, no acceleration.} Equation~\eqref{eq:cost} with
        $\lambda_a = 0$. Isolates the contribution of second-order smoothness.
  \item \textbf{B4 --- Joint, no Huber.} Squared loss in place of
        $\rho_\delta$. Isolates outlier robustification.
  \item \textbf{B5 --- World-frame.} Workspace-anchored rather than
        object-relative targets, evaluated on the displaced-object task only.
\end{enumerate}

\section{Evaluation Protocol and Metrics}
\label{sec:metrics}

Metrics were fixed before the final experimental runs, so that the criteria of
success were not selected after observing the outcomes.

\subsection{Retargeting fidelity}

Task-space tracking error between the desired end-effector pose
\eqref{eq:transfer} and the achieved pose $f(\bq_t)$, reported separately for
translation and rotation:
\begin{equation}
  \mathrm{RMSE}_{\mathrm{pos}} =
  \sqrt{\frac{1}{T}\sum_{t} \norm{\bm{p}_t^{\mathrm{des}} - \bm{p}_t}^{2}},
  \qquad
  \mathrm{RMSE}_{\mathrm{rot}} =
  \sqrt{\frac{1}{T}\sum_{t} \norm{\Log(\bm{R}_t^{\mathrm{des}\,\top}
  \bm{R}_t)}^{2}}.
  \label{eq:rmse}
\end{equation}

\subsection{Smoothness}

Spectral arc length of the joint velocity profile, a dimensionless smoothness
measure that is invariant to amplitude and duration, together with mean
absolute joint jerk. Lower magnitudes indicate smoother motion.

\subsection{Feasibility}

The proportion of synthesised episodes that execute on hardware without
controller fault, joint limit violation or collision, and the mean tracking
residual reported by the controller during replay.

\subsection{Downstream policy success}

ACT and Diffusion Policy trained on the exported dataset and evaluated over
\ph{N} rollouts per task, reporting success rate under a task-specific success
predicate.

\subsection{Calibration and sensing accuracy}

Reprojection error of the ChArUco calibration in pixels; cross-camera
registration error as the mean distance between corresponding points in
overlapping regions; and residual timestamp misalignment between device
families.

\section{Calibration and Sensing Results}
\label{sec:res-calib}

\begin{table}[htbp]
\caption{Calibration and Synchronisation Accuracy}
\label{tab:calib}
\centering
\small
\begin{tabular}{lcc}
\toprule
Quantity & Measured & Reference from literature \\
\midrule
ChArUco reprojection error (px)        & \ph{X.XX} & $<0.5$ typical \\
Cross-camera registration, IR (mm)     & \ph{XX.X} & $7.4$--$9.9$ \cite{romeo} \\
Cross-camera registration, colour (mm) & \ph{XX.X} & $20.9$--$21.4$ \cite{romeo} \\
Kinect pair timestamp drift (ms/min)   & \ph{X.XX} & $<1$ \cite{ridgerun} \\
Kinect--RealSense residual offset (ms) & \ph{X.XX} & --- \\
Wrist position repeatability (mm)      & \ph{X.XX} & --- \\
\bottomrule
\end{tabular}
\end{table}

\TODO{Populate Table~\ref{tab:calib} from the calibration runs; discuss whether
measured cross-camera registration falls within the range reported by
\cite{romeo}, and if not, whether the discrepancy is attributable to board
size, view overlap or working distance.}

\section{Retargeting Fidelity}
\label{sec:res-fidelity}

\begin{table}[htbp]
\caption{Retargeting Fidelity and Smoothness, Proposed Formulation Versus
Sequential Baseline}
\label{tab:fidelity}
\centering
\small
\begin{tabular}{lcccc}
\toprule
Method & $\mathrm{RMSE}_{\mathrm{pos}}$ (mm) & $\mathrm{RMSE}_{\mathrm{rot}}$
(deg) & SAL & Jerk (rad/s$^3$) \\
\midrule
B1 Sequential            & \ph{XX.X} & \ph{X.X} & \ph{$-$X.XX} & \ph{XX.X} \\
Proposed \eqref{eq:cost} & \ph{XX.X} & \ph{X.X} & \ph{$-$X.XX} & \ph{XX.X} \\
\midrule
Relative change          & \ph{$-$XX\%} & \ph{$-$XX\%} & \ph{$-$XX\%}
                         & \ph{$-$XX\%} \\
\bottomrule
\end{tabular}
\end{table}

\figplaceholder{45mm}
\begin{figure}[htbp]
\caption{Joint trajectory for a representative pick-and-place episode under the
sequential baseline (B1) and the proposed formulation. \TODO{insert plot}}
\label{fig:traj}
\end{figure}

\TODO{Report whether H1 is supported. State the statistical test used and the
number of episodes; if the improvement is present on one metric but not the
other, say so plainly rather than reporting only the favourable one.}

\section{Ablation Studies}
\label{sec:res-ablation}

\begin{table}[htbp]
\caption{Ablation of the Individual Terms of the Cost Functional}
\label{tab:ablation}
\centering
\small
\begin{tabular}{llcc}
\toprule
Variant & Removed term & $\mathrm{RMSE}_{\mathrm{pos}}$ (mm) & SAL \\
\midrule
Proposed & ---                        & \ph{XX.X} & \ph{$-$X.XX} \\
B2 & confidence weighting $\omega_t$   & \ph{XX.X} & \ph{$-$X.XX} \\
B3 & acceleration penalty $\lambda_a$  & \ph{XX.X} & \ph{$-$X.XX} \\
B4 & Huber robustification $\rho_\delta$ & \ph{XX.X} & \ph{$-$X.XX} \\
\bottomrule
\end{tabular}
\end{table}

\begin{table}[htbp]
\caption{Effect of Confidence Weighting Stratified by Landmark Occlusion}
\label{tab:occlusion}
\centering
\small
\begin{tabular}{lcc}
\toprule
Frame subset & B2 (uniform) & Proposed (weighted) \\
\midrule
All frames                & \ph{XX.X} & \ph{XX.X} \\
Low-occlusion frames      & \ph{XX.X} & \ph{XX.X} \\
High-occlusion frames     & \ph{XX.X} & \ph{XX.X} \\
\bottomrule
\end{tabular}

\vspace{1mm}
{\footnotesize Values are $\mathrm{RMSE}_{\mathrm{pos}}$ in mm. Occlusion
stratification uses the median of $\omega_t$ as the threshold.}
\end{table}

\TODO{Report whether H2 is supported, and specifically whether the benefit of
confidence weighting is larger on high-occlusion frames as hypothesised. If the
effect is uniform across strata, the mechanism proposed in
Section~\ref{sec:retarget} is not the operative one and this must be stated.}

\section{Replay Feasibility}
\label{sec:res-replay}

\begin{table}[htbp]
\caption{Outcome of Physical Replay Screening}
\label{tab:replay}
\centering
\small
\begin{tabular}{lccc}
\toprule
Task & Episodes synthesised & Rejected at replay & Rejection cause (dominant) \\
\midrule
Pick-and-place              & \ph{NN} & \ph{NN} (\ph{XX\%}) & \ph{[cause]} \\
Pick-and-place (displaced)  & \ph{NN} & \ph{NN} (\ph{XX\%}) & \ph{[cause]} \\
Pouring                     & \ph{NN} & \ph{NN} (\ph{XX\%}) & \ph{[cause]} \\
\midrule
Total                       & \ph{NN} & \ph{NN} (\ph{XX\%}) & --- \\
\bottomrule
\end{tabular}
\end{table}

\TODO{Report the distribution of rejection causes---joint limit, singularity
proximity, collision, controller tracking fault---since the distribution
indicates which constraint in \eqref{eq:cjoint}--\eqref{eq:ccoll} is
insufficiently modelled rather than merely how many episodes failed.}

\section{Downstream Policy Performance}
\label{sec:res-policy}

\begin{table}[htbp]
\caption{Success Rate of Policies Trained on the Exported Dataset}
\label{tab:policy}
\centering
\small
\begin{tabular}{lcc}
\toprule
Task & ACT & Diffusion Policy \\
\midrule
Pick-and-place              & \ph{XX\%} & \ph{XX\%} \\
Pick-and-place (displaced)  & \ph{XX\%} & \ph{XX\%} \\
Pouring                     & \ph{XX\%} & \ph{XX\%} \\
\midrule
Mean                        & \ph{XX\%} & \ph{XX\%} \\
\bottomrule
\end{tabular}
\end{table}

\begin{table}[htbp]
\caption{Effect of Representation Choice on the Displaced-Object Task}
\label{tab:repr}
\centering
\small
\begin{tabular}{lc}
\toprule
Representation & Success rate \\
\midrule
B5 Workspace-anchored (world frame) & \ph{XX\%} \\
Object-relative \eqref{eq:objrel}   & \ph{XX\%} \\
\bottomrule
\end{tabular}
\end{table}

\TODO{Report whether H3 is supported by comparing policy success on datasets
with and without replay screening. Note that this requires training on the
unscreened dataset as well, which doubles the training budget; if that
comparison is not run, H3 must be evaluated on the executability metric alone
and the weaker claim stated explicitly.}
